\subsection{Discussion on Evaluation and Development}

This thesis presented an exploratory investigation into building a Linux
\gls{cli} assistant powered by local \glspl{llm}. The work evaluates multiple
architectural approaches, identifying which are practical for offline,
resource-constrained environments and which prove unsuitable for reliable
real-world use. Through implementation and evaluation, significant practical insights emerged:

\paragraph{Smart Classifier Component --- Not Viable}
LLM-based classification of commands vs.\ queries performs poorly with small
local models (0.5--4B parameters). Achieving acceptable accuracy requires
either: unacceptable query latencies (2--5+ seconds per input, caused by relying
on bigger models), or extensive hardcoding the rules for classification. This
task is better handled by traditional methods and specialized models, rather
than general-purpose \glspl{llm}.

\paragraph{Agent-Based Approaches --- Unreliable In Current State} % TODO: add info in bachelor's
Agentic systems that attempt multi-step reasoning, tool use, and state management consistently
fail with small local \glspl{llm}. In current state, the approach is totaly
relying on decision-making capabilities of the model, giving it full freedom to
choose. The better approach, which could be explored in future work, is to give
agent more narrow, guided and defined set of actions, or to use model
fine-tuning to do the task.

\paragraph{RAG System --- Practical and Effective}
The retrieval-augmented generation approach proved robust and practical. By
constraining the \gls{llm} task to extracting and summarizing information from
retrieved documents, the system achieves: reliable accuracy (ranging from 80\%
to 94\% Top-5 accuracy depending on the approach), acceptable latency (under 3s
per query), and good user experience. Being the most accurate approach,
document-level retrieval has introduced the highest output size, which would
negatively affect other parts of the system, ranging around 60-100K characters
(about 50-80K tokens). To address this, passage-level retrieval and re-ranking
approaches were explored, which reduced output size to about 6K characters
(about 5K tokens) while maintaining near similar accuracy, making it more
practical for summarization and overall user experience.

\paragraph{End-to-End Latency Trade-offs}
End-to-end measurements on an NVIDIA RTX 4060 8GB GPU show that retrieval
strategy strongly affects total latency. Document-level retrieval kept search
latency around 1.58--1.60s but produced the highest RAG latency (12.39--17.21s
for the light configuration and 39.78--56.57s for the heavy configuration).
Re-ranking increased search latency (2.03--2.22s for Granite embeddings;
6.92--8.18s for qwen3-embedding) while reducing RAG latency to 5.69--7.47s
(light) and 9.57--10.77s (heavy). Passage-level retrieval kept search latency
close to document-level values (1.47--1.61s and 1.58--1.60s) while providing the
lowest RAG latency overall (5.34--6.47s light; 8.66--8.86s heavy), making it the
most practical for interactive use.

\paragraph{Practical Value}
ShAI provides practical offline assistance for \gls{cli} users, delivering fast,
help on user document database without network dependency, or privacy concerns.
The current implementation demonstrates that the \gls{rag} approach is viable
and useful, despite being less ambitious than original goals of general-purpose
command-line reasoning.

The thesis demonstrates that while ambitious visions of local \gls{llm}
assistants handling arbitrary tasks remain impractical, specialized, constrained
systems can provide real value to users with various limitations. This work lays
the groundwork for further exploration \gls{llm}-enabled information processing
and manipulation.


\subsection{Discussion of Future Directions}

\paragraph{Implications}
The findings imply that local (transformer-based) \gls{llm} assistance on
consumer hardware is best achieved through \textbf{constrained, task-specific
architectures} rather than general-purpose reasoning:
\begin{itemize}
  \item Limit task scope to information retrieval and summarization
  \item Avoid open-ended classification, reasoning, or tool-use tasks
  \item Provide focused context (via retrieval) to ground model outputs
  \item Accept trade-offs: helpful within constraints beats unreliable when
      pushed beyond them
\end{itemize}

\paragraph{Future Directions}
Future work around this topic can explore various avenues to build on the
findings of this thesis:
\begin{itemize}
  \item \textbf{Broader Task Support:} Investigate additional constrained tasks
    that can be reliably handled by local models (e.g., code generation from
    templates, error message explanations).
  \item \textbf{Application to Other Domains:} While exploring other approaches
    for small \gls{llm}-guided assistance, the focus can be shifted to other
    domains, like transport assistants, personal information management,
    mobile \gls{gui} apps, etc.
  \item \textbf{User Interface Improvements:} Develop richer interfaces (e.g.,
    interactive query refinement, multi-turn conversations) while maintaining
    low latency and simplicity. Also further exploration of different seamless
    integration approaches, used in solutions like \texttt{llm-zsh-plugin}
    or \texttt{AiTerm}, can be done.
  \item \textbf{Software Optimization/Refactoring:} Explore the possibilities of
    usage of LangChain, or other similar frameworks, adapt existing software
    for better performance (models caching), implement MCP-based architecture,
    tools, optimize architecture for being a framework for building small
    automations for everyday \gls{cli} usage.
  \item \textbf{Experimenting with top\_k and section size for passage-level
      retrieval:} The thesis did not describe the exploration of different
        parameters for passage-level retrieval, as the goal of achieving the
        acceptable accuracy with the lowest latency was achieved, but further
        exploration of different parameters for passage-level retrieval can be
        done, with the software already implemented.
\end{itemize}
Additionally, if we look at the future directions from the perspective not
limited to the scope of the developed application, but more generally, we can
define further research directions in the field:
\begin{itemize}
  \item \textbf{Comparing Different Model Testing Approaches and Frameworks:}
    The thesis does not go deep into the comparison of different approaches for
    testing and evaluating \gls{llm}-based systems, but this is very
    important area for the research, as it provides necessary empirical
    evidence that proves the results of any research.
  \item \textbf{Model Optimization and Architectures Exploration:} Explore model
    distillation, quantization, or fine-tuning to further improve performance
    and reduce latency on local hardware. With the rapid evolution of
    \gls{llm}, new architectures like \gls{rnn}-augmented \glspl{llm},
    Mamba, \gls{lstm}-based models, etc. can be explored for their potential
    in constrained tasks. Future work can focus around testing and
    evaluating these models for constrained environments, but also be more
    focused on the development of new architectures that are optimized for
    constrained environments and tasks.
  \item \textbf{Caching and Context Management:} A lot of tools nowadays support
    as-called ``context snapshots'', which are basically a way to save the state
    of the system at a given point in time. This is useful when working with
    repeating situations that need previous setup. Caching can potentially
    have huge impact on the performance of the system, saving time and
    resources.
\item \textbf{Exploring Different Harnessing Approaches:} Apart from \gls{rag}
    approach, there are various other ways to use the \gls{llm} capabilities in
    different ways, including but not limited to: different agentic
    approaches (for example, Program-Aided Language Models (PAL) /
    Program-of-Thought (PoT)), different classification approaches (e.g.,
    small specialized models, not limited to \gls{llm}).
\end{itemize}


\subsection{Contributions}

This work contributes to understanding practical \gls{llm} integration in
terminal environments:

\begin{itemize}
  \item Empirical evidence of which architectures work with small local \glspl{llm}
  \item The software for evaluating the approaches for \gls{rag} and
      classification, that can be used for further research in this area, in
      scope of accuracy, latency, and user experience of \gls{llm}-based
      systems.
  \item Identification of key failure modes in applying local models to structured tasks
  \item An open-source reference implementation of a working \gls{rag}-based
      assistant, which can also be used as a base framework for further
      improvements and automations in scope of local \gls{llm} assistance in
      terminal environments.
  \item Guidance for future research: local model assistance is viable within constrained,
      well-scoped tasks, not as general-purpose reasoning
\end{itemize}
